% known_bad fixture (D-20261012-3, Louis round 3): offending sentences of the 17:35 tree, verbatim.

\section{Fixture}

Table~\ref{tab:coverage} presents coverage probabilities of the 95\% bootstrap confidence intervals. For the non-null scenarios S3--S7 at $n \geq 100$, coverage ranged over \numval{covRangeS3toS7N100up2dp}: S3 (\numval{covS3rangeN100up}), S4 (\numval{covS4rangeN100up}), S5 (\numval{covS5n100}--\numval{covS5n200}), S6 (\numval{covS6rangeN100up}), and S7 (\numval{covS7rangeN100up}); of these ten cells, \numval{nCellsCovBelow2MCSE} lie more than two Monte Carlo standard errors below 0.95, the lowest being S5 at $n = 100$ (\numval{covS5n100}). The shortfall is largest in S5 and S7, whose CDFs cross; there the intervals lie above the true value, in line with the upward bias of Section~\ref{sec:estimation}, which relative to the sampling standard deviation of $\widehat{W}_1$ is \numval{biasOverSdS5n100} and \numval{biasOverSdS7n100} at $n = 100$ against \numval{biasOverSdS3n100} for S3. In S6 (\numval{biasOverSdS6n100}) and in S4 at $n = 100$, whose CDFs do not cross and whose bias is within Monte Carlo error, the shortfall is small and the misses fall on both sides (Table~\ref{tab:study2_intervals}), so a given $n$ does not by itself fix the coverage.

Monte Carlo standard errors for the reported coverage probabilities are at most \numval{mcseBoundPropB} at 10{,}000 replications (largest observed \numval{mcseCovMax}). A difference between two cells is within simulation noise when it is smaller than about $\numval{zDiffFactor}$ times the standard error, and a single cell is consistent with nominal coverage when it lies within about twice its own standard error of 0.95. Because the application reads an unresolved distance at its upper limit, the one-sided coverage of the 97.5\% upper limit, $P(W_1 \leq U)$, is the operating characteristic that matters there: for S3--S7 at $n \geq 100$ it ranged over \numval{covUpperMinS3toS7N100up}--\numval{covUpperMaxS3toS7N100up} against a nominal 0.975, so the upper limit is conservative wherever the interval undercovers. The basic (reflected) interval $[\max(0, 2\widehat{W}_1 - U), \max(0, 2\widehat{W}_1 - L)]$, which shifts the percentile interval by the bootstrap estimate of the bias, covered less than the percentile interval in every non-null cell (\numval{covBasicRangeS3toS7N100up} at $n \geq 100$; under S1 its coverage, \numval{covBasicS1n50} at $n = 50$, equals the share of replications whose lower limit is truncated at zero), so bias correction of this kind does not repair the boundary behaviour (Table~\ref{tab:study2_intervals}); these comparisons were computed after the simulation from the stored replications and were not pre-specified.

Monte Carlo standard errors for the reported coverage probabilities are at most \numval{mcseBoundPropB} at 10{,}000 replications (largest observed \numval{mcseCovMax}). A difference between two cells is within simulation noise when it is smaller than about $\numval{zDiffFactor}$ times the standard error, and a single cell is consistent with nominal coverage when it lies within about twice its own standard error of 0.95. Because the application reads an unresolved distance at its upper limit, the one-sided coverage of the 97.5\% upper limit, $P(W_1 \leq U)$, is the operating characteristic that matters there: for S3--S7 at $n \geq 100$ it ranged over \numval{covUpperMinS3toS7N100up}--\numval{covUpperMaxS3toS7N100up} against a nominal 0.975, so the upper limit is conservative wherever the interval undercovers. The basic (reflected) interval $[\max(0, 2\widehat{W}_1 - U), \max(0, 2\widehat{W}_1 - L)]$, which shifts the percentile interval by the bootstrap estimate of the bias, covered less than the percentile interval in every non-null cell (\numval{covBasicRangeS3toS7N100up} at $n \geq 100$; under S1 its coverage, \numval{covBasicS1n50} at $n = 50$, equals the share of replications whose lower limit is truncated at zero), so bias correction of this kind does not repair the boundary behaviour (Table~\ref{tab:study2_intervals}); these comparisons were computed after the simulation from the stored replications and were not pre-specified.

Table~\ref{tab:study2_required_n_mc} reports, for all \numval{nCellsSelection} cells, the per-region sample size at which each measure's selection AUC reaches 0.9, with 95\% Monte Carlo intervals. Two kinds of number appear in it, and they answer different questions. Reading down the $W_1$ column shows how $W_1$'s requirement varies with the task, from \numval{reqnW1Set4symsev} (symmetric severity) to \numval{reqnW1maxSets34b} (shape-symmetric mixtures) in Sets 3 and 4. Reading across a row gives the method difference within one task: in the symmetric-severity cell, $W_1$ reaches the target at $n = \numval{reqnW1Set4symsevB}$ and the KS statistic at $n = \numval{reqnKSSet4symsev}$, a \numval{reqnRatioKSoverW1Set4symsev}-fold difference on the same data-generating process. At the AUC 0.9 target the 95\% intervals place $W_1$ below the KS statistic in \numval{nCellsW1FewerKS} cells---three of them the Set 4 subgroup cells, where the advantage follows from the construction (Section~\ref{sec:study1_design}), and among the others all six light-tailed cells of Sets 1 and 2, where neither measure is favoured by construction---place the KS statistic below $W_1$ in \numval{nCellsKSFewer}, the bulk-shift control, and overlap in \numval{nCellsOverlap}, the shape-symmetric mixture (\numval{reqnW1Set3shapesym} [\numval{reqnW1Set3shapesymLo}, \numval{reqnW1Set3shapesymHi}] against \numval{reqnKSSet3shapesym} [\numval{reqnKSSet3shapesymLo}, \numval{reqnKSSet3shapesymHi}]). In the bulk-shift control the KS statistic reaches the target first (\numval{reqnKSSet4bulk} versus \numval{reqnW1Set4bulk}), whereas in the light-tailed location cells of Sets 1 and 2 $W_1$ does (\numval{reqnW1Set1loc} versus \numval{reqnKSSet1loc}; \numval{reqnW1Set2loc} versus \numval{reqnKSSet2loc}). A registered check (protocol and outputs in the repository) locates the cause in the sampling variability of $\widehat{W}_1$: for the same two-unit shift at $n = 200$ per region, the Set 4 anchor population, whose two 5\% subgroups lie 40 units from the centre, multiplies the null sampling standard deviation of $\widehat{W}_1$ by \numval{set4NullSdRatioW1n200} (Monte Carlo SE \numval{set4NullSdRatioW1n200Mcse}) relative to $N(50, 10^2)$ while leaving that of the distribution-free KS statistic unchanged (\numval{set4NullSdRatioKSn200}, SE \numval{set4NullSdRatioKSn200Mcse}), and the paired AUC difference between $W_1$ and the KS statistic changes sign, from \numval{set4DeltaAucGaussN200d2} (SE \numval{set4DeltaAucGaussN200d2Mcse}) under the Gaussian anchor to \numval{set4DeltaAucMixN200d2} (SE \numval{set4DeltaAucMixN200d2Mcse}) under the Set 4 anchor; in a population with such contaminated tails a location difference favours the KS statistic, and we report the loss.

Table~\ref{tab:study2_required_n_mc} reports, for all \numval{nCellsSelection} cells, the per-region sample size at which each measure's selection AUC reaches 0.9, with 95\% Monte Carlo intervals. Two kinds of number appear in it, and they answer different questions. Reading down the $W_1$ column shows how $W_1$'s requirement varies with the task, from \numval{reqnW1Set4symsev} (symmetric severity) to \numval{reqnW1maxSets34b} (shape-symmetric mixtures) in Sets 3 and 4. Reading across a row gives the method difference within one task: in the symmetric-severity cell, $W_1$ reaches the target at $n = \numval{reqnW1Set4symsevB}$ and the KS statistic at $n = \numval{reqnKSSet4symsev}$, a \numval{reqnRatioKSoverW1Set4symsev}-fold difference on the same data-generating process. At the AUC 0.9 target the 95\% intervals place $W_1$ below the KS statistic in \numval{nCellsW1FewerKS} cells---three of them the Set 4 subgroup cells, where the advantage follows from the construction (Section~\ref{sec:study1_design}), and among the others all six light-tailed cells of Sets 1 and 2, where neither measure is favoured by construction---place the KS statistic below $W_1$ in \numval{nCellsKSFewer}, the bulk-shift control, and overlap in \numval{nCellsOverlap}, the shape-symmetric mixture (\numval{reqnW1Set3shapesym} [\numval{reqnW1Set3shapesymLo}, \numval{reqnW1Set3shapesymHi}] against \numval{reqnKSSet3shapesym} [\numval{reqnKSSet3shapesymLo}, \numval{reqnKSSet3shapesymHi}]). In the bulk-shift control the KS statistic reaches the target first (\numval{reqnKSSet4bulk} versus \numval{reqnW1Set4bulk}), whereas in the light-tailed location cells of Sets 1 and 2 $W_1$ does (\numval{reqnW1Set1loc} versus \numval{reqnKSSet1loc}; \numval{reqnW1Set2loc} versus \numval{reqnKSSet2loc}). A registered check (protocol and outputs in the repository) locates the cause in the sampling variability of $\widehat{W}_1$: for the same two-unit shift at $n = 200$ per region, the Set 4 anchor population, whose two 5\% subgroups lie 40 units from the centre, multiplies the null sampling standard deviation of $\widehat{W}_1$ by \numval{set4NullSdRatioW1n200} (Monte Carlo SE \numval{set4NullSdRatioW1n200Mcse}) relative to $N(50, 10^2)$ while leaving that of the distribution-free KS statistic unchanged (\numval{set4NullSdRatioKSn200}, SE \numval{set4NullSdRatioKSn200Mcse}), and the paired AUC difference between $W_1$ and the KS statistic changes sign, from \numval{set4DeltaAucGaussN200d2} (SE \numval{set4DeltaAucGaussN200d2Mcse}) under the Gaussian anchor to \numval{set4DeltaAucMixN200d2} (SE \numval{set4DeltaAucMixN200d2Mcse}) under the Set 4 anchor; in a population with such contaminated tails a location difference favours the KS statistic, and we report the loss.

The sponsor selects Region~8 as the anchor region---representing, for instance, a regulatory market where local efficacy evidence is required. Region~8 comprises $n = \numval{nR8}$ patients in GUSTO-I, with baseline characteristics summarized in Table~\ref{tab:gusto_r8}; the anchor role, not its size, is what the exercise illustrates. Quantifying how each region differs from Region~8 mirrors the input a sponsor would bring to a pooling judgment as described in ICH E17.\cite{iche17} The GUSTO-I records are used as recorded: \numval{nSbpZero} patients have a systolic blood pressure of 0~mmHg and \numval{nAgeOver100} are older than 100 years, and $W_1$ weights each observation by its displacement. Excluding these records changes no $\widehat{W}_1$ by more than \numval{w1MaxChangeExclAge} years or \numval{w1MaxChangeExclSbp} mmHg, exchanges \numval{nRankChangesExcl} adjacent SBP ranks and changes no null-floor classification below; winsorizing both variables at the 1st and 99th percentiles moves $\widehat{W}_{1,\text{SBP}}$ by up to \numval{w1MaxChangeWinsorSbp} mmHg. Systolic blood pressure is also heaped: the share of readings on a multiple of 10~mmHg ranges from \numval{sbpDigit0Min}\% to \numval{sbpDigit0Max}\% across the 16 regions and is \numval{sbpDigit0R8}\% in Region~8, so differences in recording practice contribute to the SBP distances and the step CDFs of Appendix~\ref{app:asymptotic} apply.

\abstract[Abstract]{%
The ICH E17 guideline describes regional pooling in multi-regional clinical trials (MRCTs) in terms of the similarity of effect modifier distributions but names no measure, and to our knowledge no distance between regional effect modifier distributions has been evaluated for it. We quantify how two regions' distributions of each candidate effect modifier differ by the Wasserstein-1 ($W_1$) distance with bootstrap confidence intervals. If the conditional average treatment effect is common to all regions, depends on that modifier alone, and is Lipschitz with clinical slope $L_{\text{clinical}}$, regional average effects differ by at most $\Delta_{\max} = L_{\text{clinical}} W_1$; when $L_{\text{clinical}}$ is unknown, the smallest slope $L^*$ at which the distance could produce a clinical margin is compared with a sponsor's bound. In a selection simulation with the number of matches supplied, the standardized mean difference stayed at chance up to 2,000 patients per region in some configurations; $W_1$ and the Kolmogorov--Smirnov (KS) statistic reached the target (AUC 0.9) in every configuration, $W_1$ needing fewer patients in \numval{nCellsW1FewerKS} of \numval{nCellsSelection} cells (three by construction), overlapping in one, and more where a location shift sat on contaminated tails. For true distances of about \numval{minW1SdS3toS7} to \numval{maxW1SdS3toS7} SD (normal, equal $n$), bootstrap coverage from 100 patients per region was \numval{covRangeS3toS7N100up2dpAbs}, below 0.95 beyond Monte Carlo error in \numval{nCellsCovBelow2MCSE} of ten cells. In a GUSTO-I-based hypothetical MRCT, partner rankings on age and on systolic blood pressure differed; a null floor marked distances not distinguishable from no difference. The framework quantifies differences; the pooling judgment remains clinical.}
